\chapter{Proposed Modeling Approaches}
\label{chap:task-formulations}

The modeling strategy is organized around four predictive formulations rather
than around standalone architectures. All four approaches share the same
operational objective: at decision time \(t_i\), use only the signal history
available up to that time to decide whether the backup link should be activated.
They differ in the intermediate quantity estimated before this final decision:
the future signal trajectory, a local persistence duration, the probability of
a long fade event, or a residual-duration survival function.

\section{Overview of the Four Modeling Formulations}
\label{sec:overview-modeling-formulations}

The four formulations represent different levels of abstraction of the same
nowcasting problem. Autoregressive forecasting retains the most detailed
description of the future by predicting a trajectory. Current-level persistence
reduces the future to the time for which the signal remains close to its current
level. Long-fade detection formulates the problem as an event-level binary
classification task. Survival persistence retains uncertainty about the time to
recovery by estimating a residual-duration distribution while allowing for
right-censored observations.

This organization separates the predictive task from the architecture used to
implement it. The following matrix summarizes the target, the model families
used in the corresponding experiments, and the conversion that connects each
native output to the common switch decision. The detailed definitions and
assumptions are given in \cref{sec:autoregressive-formulation,sec:current-level-formulation,sec:long-fade-formulation,sec:survival-formulation}.

\begin{table}[htbp]
  \centering
  \small
  \caption{Overview of the four modeling formulations and their implemented
  model families.}
  \label{tab:modeling-formulation-matrix}
  \vspace{0.45em}
  \renewcommand{\arraystretch}{1.35}
  \setlength{\tabcolsep}{5pt}
  \begin{tabular}{@{}>{\raggedright\arraybackslash}p{0.16\textwidth}
    >{\raggedright\arraybackslash}p{0.21\textwidth}
    >{\raggedright\arraybackslash}p{0.30\textwidth}
    >{\raggedright\arraybackslash}p{0.18\textwidth}@{}}
    \toprule
    Formulation & Native output & Model families & Switch conversion \\
    \midrule
    Autoregressive forecasting (Sec.~\ref{sec:autoregressive-formulation})
      & Future signal trajectory
      & GRU sequence models \cite{cho2014learning,chung2014empirical};
        PatchTST \cite{nie2023patchtst}; XGBoost \cite{chen2016xgboost};
        Chronos \cite{ansari2024chronos}
      & Forecast trajectory to switch \\
    Current-level persistence (Sec.~\ref{sec:current-level-formulation})
      & Local persistence duration
      & Multiscale learnable-shapelet MLP, Transformer, and convolution
        variants \cite{grabocka2014learning}; XGBoost \cite{chen2016xgboost}
      & Duration threshold \\
    Long-fade detection (Sec.~\ref{sec:long-fade-formulation})
      & Probability of a long grouped fade event
      & XGBoost lag-and-scalar classifier \cite{chen2016xgboost}; TCN
        classifier \cite{bai2018tcn}; multiscale shapelet convolution
        classifier \cite{grabocka2014learning}
      & Probability threshold with signal gate \\
    Survival persistence (Sec.~\ref{sec:survival-formulation})
      & Residual-duration survival function
      & XGBoost-AFT \cite{barnwal2022xgboostAft}; discrete-time causal TCN
        \cite{gensheimer2019nnet,bai2018tcn}
      & Survival probability at the operational horizon \\
    \bottomrule
  \end{tabular}
\end{table}

The next sections define the common prediction setting and then describe each
formulation in detail. For each one, the input representation, target
construction, model families, native metrics, and output-to-switch rule are
specified before the formulations are compared in a common operational
evaluation framework.

\section{Common Prediction Setting}
\label{sec:common-prediction-setting}

After the common loading and preprocessing stage, each dataset is represented as
an ordered univariate time series

\[
  \{(t_i, S_i)\}_{i=1}^{N},
\]

where \(t_i\) is the timestamp of the \(i\)-th observation and \(S_i\) is the
corresponding signal value. The project deliberately uses a signal-only
representation: no meteorological, orbital, station-side, or file-specific
auxiliary variable is used as model input. This constraint makes the comparison
between datasets and tasks consistent, because every model receives information
derived from the same physical quantity.

For a decision made at index \(i\), the generic past context of length \(L\) is

\[
  X_i = [S_{i-L+1}, \ldots, S_i].
\]

The last value \(S_i\) is the most recent observation available when the
decision is made. All features used by the models must be observable at time
\(t_i\): they can be functions of this past context, or task-specific
elapsed-time and event-state quantities computed from past and current
threshold crossings. Future samples may be used to construct supervised targets
during dataset generation, but they are never allowed to enter \(X_i\) or any
derived input feature.

Different tasks apply different transformations to the same basic context. Some
models use the raw window directly, some use a local context-standardized
version, and some use relative representations such as \(X_i-S_i\) or
\(X_i-\theta\), where \(\theta\) is the operational signal threshold. When
scalar context features are used, they are computed only from the available
signal history and online task state, and their standardization parameters are
estimated on the training split only.

The term scalar context is used in a precise sense throughout the chapter. It
does not denote external covariates. It denotes past-only scalar summaries
derived from the same signal window and, when the task is event-based, from the
event state already observable at \(t_i\). The recurring signal summaries are
trailing quantities:
\[
  \bar{S}_{i,L}
  =
  \frac{1}{L}\sum_{r=0}^{L-1} S_{i-r},
  \qquad
  \operatorname{sd}_{i,L}
  =
  \operatorname{std}(S_{i-L+1},\ldots,S_i),
\]
together with the window minimum, maximum, range, and finite differences
\[
  \Delta_i^{(q)} = S_i - S_{i-q+1},
  \qquad q\in\{2,3,5\}.
\]
The short-window summaries \(\bar{S}_{i,5}\) and
\(\operatorname{sd}_{i,5}\) are computed on the last five available samples of
the context. These are rolling only in the trailing online sense: they are never
centered around \(t_i\), and they never use samples after \(t_i\). When a model
uses a tabular lag matrix, the lag columns are ordered from oldest to newest;
\(\operatorname{lag}_{L-1}\) corresponds to \(S_{i-L+1}\) after the selected
task transformation, and \(\operatorname{lag}_{0}\) corresponds to the current
sample \(S_i\) after the same transformation.

The output of a model at index \(i\) is denoted generically as \(z_i\). Depending
on the task, \(z_i\) can be a future trajectory, a duration estimate, a binary
event probability, or a survival curve. Each task therefore defines an explicit
conversion rule

\[
  r_i = c(z_i, S_i),
\]

where \(r_i \in \{0,1\}\) is the raw model-derived switch before the shared
post-processing defined in \cref{sec:shared-switch-post-processing}. The final
evaluated switch is obtained only after applying that common policy. This
separation keeps the task-specific predictive problem distinct from the common
operational decision layer.

\section{Autoregressive Forecasting}
\label{sec:autoregressive-formulation}

Autoregressive forecasting is the most direct formulation of the nowcasting
problem considered in this work. Instead of predicting a binary switch or a
duration variable directly, the model predicts the future signal trajectory over
a short operational horizon. The switch decision is then derived from this
forecasted trajectory. This formulation is useful because it preserves an
interpretable intermediate output: before evaluating whether a method triggers a
backup-link switch, one can inspect whether it predicts the local evolution of
the signal correctly.

This signal-forecasting view deliberately avoids building a fade definition
into the predictive target. The model is not trained to decide whether an event
is a fade, whether a fade will last longer than a fixed duration, or whether a
particular threshold crossing belongs to an operationally relevant event. It is
trained only to estimate future values of the measured signal. As a consequence,
the autoregressive task does not require a threshold \(\theta\), an onset or
recovery rule, an event-grouping rule, or a minimum fade duration in order to be
defined. These quantities become necessary only later, when the predicted signal
trajectory is converted into a switch decision and compared with reference
switch policies. This is a key difference from the task formulations that
follow, where the supervised target itself is a duration, event class, or
survival quantity whose definition depends on a fade or recovery criterion.

\subsection{Task Definition}

For a valid decision index \(i\), the autoregressive input is the past signal
window

\[
  X_i^{\text{AR}} = [S_{i-L+1}, \ldots, S_i] \in \mathbb{R}^{L},
\]

and the target is the future signal trajectory

\[
  Y_i^{\text{AR}} = [S_{i+1}, \ldots, S_{i+h}] \in \mathbb{R}^{h}.
\]

Both quantities are defined only from the signal and the forecasting horizon.
A valid sample requires an observed context and a fully observed future horizon
inside one continuous acquisition segment. In the experiments, the eligible
decision indices are restricted to the candidate event windows built during data
preparation. This is a sampling and evaluation-design choice: it focuses the
training and comparison on operationally relevant periods, but it does not
change the autoregressive target, which remains the future signal trajectory.

In the main experimental configuration, \(L=30\) and \(h=10\). With the
canonical sampling interval of approximately 30 seconds, the model therefore
observes about 15 minutes of recent signal history and predicts about 5 minutes
ahead. A longer-context Chronos diagnostic configuration with \(L=120\) is also
kept separate from the main autoregressive comparison.

The future trajectory is used only to construct the offline target. The model
output is

\[
  \widehat{Y}_i^{\text{AR}}
  =
  [\widehat{S}_{i+1}, \ldots, \widehat{S}_{i+h}],
\]

and is converted back to the raw signal scale before forecast metrics or switch
conversion are computed.

\subsection{Tested Models}

Four autoregressive model families are used. Their general modeling principles
are introduced in \cref{chap:background}; this section only specifies how they
are instantiated for the autoregressive task. For trainable neural models, the
input tensor has shape \(N\times L\times 1\), where the last dimension is the
single \texttt{Signal} channel. Two input variants are considered. In the raw
variant, \(X_i^{\text{AR}}\) and \(Y_i^{\text{AR}}\) are used in their original
scale. In the context-standardized variant, each context is normalized with
statistics computed from the context itself, and the corresponding target is
represented in the same local scale during training. Predictions from this
variant are always inverse-transformed to raw signal units before forecast
metrics or switch decisions are computed.

This two-variant design is used for every trainable autoregressive model family
in the experiments: GRU, PatchTST, and XGBoost. Chronos is evaluated separately
as a raw zero-shot foundation-model baseline, because it performs its own
internal scaling.

For the context-standardized variant, the dataset stores
\[
  \mu_i=\frac{1}{L}\sum_{r=0}^{L-1}S_{i-r},
  \qquad
  s_{\mathrm{ctx},i}=\max\{\operatorname{std}(X_i^{\text{AR}}),\epsilon\}.
\]
Here \(s_{\mathrm{ctx},i}\) is the context standard deviation bounded below by
\(\epsilon>0\) for numerical stability. No rolling statistics, slope features,
scalar context vector, threshold mask, or event-position variable is appended to
the autoregressive models. GRU, PatchTST, and Chronos receive only the ordered
univariate sequence. The XGBoost autoregressive baseline receives the same values
flattened into lag columns \(\operatorname{lag}_{L-1},\ldots,\operatorname{lag}_0\),
but no additional summary features are added.

\paragraph{GRU sequence models.}
Two GRU variants with unidirectional encoders are evaluated. Unidirectionality
is an architectural choice in these experiments; as discussed in
\cref{sec:background-gru}, a bidirectional encoder restricted to the same past
context would also respect the decision-time information boundary. Let
\(\widetilde{X}_i\) denote the model input in either raw or
context-standardized scale. The sequence-to-vector variant encodes the full
past context with a unidirectional GRU,
\[
  (A_{i,1},\ldots,A_{i,L}), a_i
  =
  \operatorname{GRU}_{\mathrm{enc}}(\widetilde{X}_i),
\]
where \(a_i\) is the final hidden representation of the last recurrent layer.
The complete future trajectory is then produced by one linear output head:
\[
  \widehat{Y}_i^{\text{AR}} = W_o a_i + b_o .
\]
This is a direct multiple-output forecaster. It shares the encoder across the
context but does not generate the forecast recursively; all horizons are
predicted from the same encoded state.

The encoder-decoder variant uses the encoder hidden state to initialize a second
GRU that runs over the future horizon. The decoder starts from the last observed
input value and generates one value at a time:
\[
  d_{i,k}
  =
  \operatorname{GRU}_{\mathrm{dec}}(\xi_{i,k},d_{i,k-1}),
  \qquad
  \widehat{S}_{i+k}=W_d d_{i,k}+b_d .
\]
During training, the decoder input \(\xi_{i,k}\) can be either the previous
prediction or the corresponding previous true target value according to the
configured teacher-forcing ratio. During validation, testing, and switch
conversion, teacher forcing is disabled, so future true values never enter the
decoded trajectory.

Both GRU variants are optimized with mean squared error over the \(h\)-step
target,
\[
  \mathcal{L}_{\mathrm{AR}}
  =
  \frac{1}{Nh}
  \sum_i\sum_{k=1}^{h}
  \left(\widehat{S}_{i+k}-S_{i+k}\right)^2 ,
\]
computed in the training scale of the selected variant.

\paragraph{PatchTST.}
PatchTST receives the same univariate context, partitions it into local temporal
patches, processes the resulting tokens with Transformer encoder blocks, and
maps the encoded representation to the \(h\)-step forecast. In this formulation,
PatchTST tests whether patch-level attention is useful for representing short
fade rises, plateaus, and recoveries within the recent signal history. The
implemented model is a deterministic PatchTST-style direct forecaster rather
than a full NeuralForecast wrapper: it extracts complete patches from the
single-channel signal context, applies a learned linear patch embedding and a
learned positional embedding, encodes the patch sequence with a Transformer
encoder, flattens the encoded tokens, and projects them to the full future
signal trajectory. Channel independence is therefore satisfied trivially,
because the task uses only the \texttt{Signal} input channel.

More explicitly, for patch length \(P\) and stride \(s\), the input is unfolded
into \(N_p\) complete patches \(p_{i,j}\in\mathbb{R}^{P}\). Each patch is mapped
to a token
\[
  e_{i,j}=W_p p_{i,j}+b_p+r_j ,
\]
where \(r_j\) is a learned positional embedding. The Transformer encoder returns
contextualized patch tokens \(H_i\), and the output head applies a direct
multi-output projection:
\[
  \widehat{Y}_i^{\text{AR}}
  =
  W_o\operatorname{vec}(H_i)+b_o .
\]
The model is trained with the same trajectory MSE used for the GRU forecasters.
Unlike the GRU encoder-decoder, it does not autoregress over the predicted
future values; the whole horizon is produced in one forward pass.

\paragraph{XGBoost autoregressive baseline.}
The XGBoost baseline converts each context into a tabular lag vector
\[
  \phi_i^{\text{AR}}
  =
  [\widetilde{S}_{i-L+1},\ldots,\widetilde{S}_i],
\]
again using either the raw or context-standardized variant. It then trains one
independent boosted-tree regressor for each forecast horizon:
\[
  \widehat{S}_{i+k}=F_k(\phi_i^{\text{AR}}),
  \qquad k=1,\ldots,h .
\]
Each \(F_k\) is an XGBoost regressor with a squared-error regression objective.
The horizon-specific regressors have separate parameters and random seeds, so
the method does not impose a joint trajectory structure across the ten predicted
future samples. Its purpose is to test whether a strong nonlinear learner over
fixed lag features can compete with explicit sequence models for short-horizon
signal prediction. Because the features are fixed lags, the baseline is causal
by construction as long as \(\phi_i^{\text{AR}}\) is built only from the past
context.

\paragraph{Chronos foundation model: zero-shot forecasting.}
Chronos is used only in zero-shot mode. The pretrained model receives a
univariate context and returns sampled future trajectories
\[
  \widehat{Y}_{i}^{(b)}
  \sim
  p_{\mathrm{Chronos}}(Y_i\mid X_i),
  \qquad b=1,\ldots,B .
\]
The deterministic trajectory used for MAE, RMSE, and switch conversion is the
sample median at each horizon,
\[
  \widehat{S}_{i+k}
  =
  \operatorname{median}_{b=1,\ldots,B}
  \widehat{S}_{i+k}^{(b)} .
\]
Prediction quantiles are also saved as distributional diagnostics, but the
switch pipeline uses the median trajectory for consistency with the
deterministic forecasters. No Chronos parameter is trained or fine-tuned, and no
train or validation data are loaded to adapt the model. A main run uses
\(L=30\), while an additional diagnostic run uses \(L=120\) to test whether a
longer context improves the foundation-model forecast.

\subsection{Switch Conversion}

The autoregressive model does not output a switch directly, and the threshold
used below is not a training label. It belongs to the downstream decision layer:
after the future signal trajectory has been predicted, the forecast is converted
into one raw switch decision at the prediction time \(t_i\). Let \(\theta\) be
the operational signal threshold and let \(q\) be the configured minimum number
of forecasted future samples that must exceed the threshold. The raw
autoregressive switch is

\[
  r_i^{\text{AR}}
  =
  \mathbb{1}
  \left[
    \sum_{k=1}^{h}
    \mathbb{1}\left(\widehat{S}_{i+k} \ge \theta\right)
    \ge q
  \right].
\]

In the main experiments, \(\theta=10.0\) dB and \(q=h=10\). Therefore the model
switches at time \(t_i\) only if all ten predicted future signal values are at
or above the operational threshold. This is a deliberately conservative rule:
it requires the forecast to indicate persistence throughout the full prediction
horizon before the backup-link decision is activated.

Changing \(\theta\), \(q\), or the post-processing rule would change the
operational interpretation of the same predicted trajectory, but it would not
change the autoregressive supervised learning problem. This separation is useful
when the desired fade definition is uncertain or may be revised, because the
forecast remains a signal-level object rather than a model-specific fade label.

After this raw decision is produced, the shared switch post-processing defined
in \cref{sec:shared-switch-post-processing} is applied. The autoregressive
models are therefore compared against the Perfect Switch using the same final
decision layer as the other task formulations. Forecast MAE and RMSE remain
useful diagnostics, but the operational comparison is based on the derived
switch sequence.

\section{Current-Level Persistence}
\label{sec:current-level-formulation}

Current-level persistence reformulates the problem as a duration regression
task. Instead of predicting the future signal trajectory, the model estimates
how long the signal will remain above a level defined relative to the current
observation. This is not equivalent to forecasting the end of the whole fade
event. It asks a local question: given the recent signal pattern and the current
level \(S_i\), for how long does the signal persist above \(S_i-\delta\)?

This formulation was included because the operational switch decision is
ultimately related to persistence. A short threshold crossing may not justify
switching to a backup link, while a signal condition expected to persist for
several minutes may. Current-level persistence therefore compresses the future
trajectory into a single scalar duration.

\subsection{Task Definition}

For a valid decision index \(i\), the raw context is again

\[
  X_i = [S_{i-L+1}, \ldots, S_i].
\]

The task-specific recovery level is defined as

\[
  \lambda_i = S_i - \delta,
\]

where \(\delta>0\) controls how far below the current value the signal is
allowed to fall before the current-level persistence episode is considered to
have ended. In the experiments, \(\delta=0.5\) dB. The target recovery index is

\[
  \tau_i =
  \inf \left\{
    j>i : S_j < \lambda_i
  \right\},
\]

provided that such an index is observed before the end of the current
acquisition segment. The observed remaining persistence duration is

\[
  D_i^{\text{CLP}} = t_{\tau_i} - t_i .
\]

If no recovery is observed before the acquisition segment ends, the candidate is
right-censored. In the current duration-regression implementation these
censored candidates are stored for diagnostics but excluded from ordinary
supervised regression. They are not assigned an artificial finite duration,
because doing so would bias the target downward.

The primary sequence representation is relative to the current value:

\[
  X_i^{\text{rel-current}}
  =
  X_i - S_i
  =
  [S_{i-L+1}-S_i, \ldots, 0].
\]

This removes the absolute level from the sequence and emphasizes local shape:
rising history, recovery history, recent oscillations, and short-term
stability. Absolute level information is reintroduced through scalar context
features computed from \(X_i\), such as the current signal, the recovery level,
window summary statistics, recent slopes, and short-window moments. These
features are standardized using training-split statistics only.

In the implemented current-level dataset, the sequence input is the
current-relative array, and the scalar vector is the scalar-context feature
array. For one sample, the scalar vector is ordered as
\[
\begin{aligned}
  s_i^{\text{CLP}} = [
  &S_i,\ S_i-\delta,\ \bar{S}_{i,L},\
  \operatorname{sd}_{i,L},\
  \min(X_i),\ \max(X_i),\
  \max(X_i)-\min(X_i),\\
  &\Delta_i^{(2)},\ \Delta_i^{(3)},\ \Delta_i^{(5)},\
  \bar{S}_{i,5},\ \operatorname{sd}_{i,5}
  ] .
\end{aligned}
\]
Thus the learnable sequence component sees only deviations from the current
sample, while the scalar branch receives the absolute level, the recovery level
\(S_i-\delta\), full-window trailing summaries, and short-window trailing
summaries. No fade-threshold variable, event-duration variable, future recovery
time, or target duration is included as an input feature.

The regression target is log-compressed:

\[
  y_i^{\text{CLP}} = \log(1+D_i^{\text{CLP}}).
\]

Models are trained to predict \(y_i^{\text{CLP}}\). For evaluation and switch
conversion, predictions are mapped back to seconds:

\[
  \widehat{D}_i^{\text{CLP}}
  =
  \exp(\widehat{y}_i^{\text{CLP}})-1 .
\]

The logarithmic target reduces the influence of very long persistence episodes
while preserving the ordering of durations.

\subsection{Tested Models}

The current-level persistence branch uses learnable shapelet models and an
XGBoost duration-regression baseline. Shapelet-based modeling is introduced in
\cref{chap:background}; here the important point is that the learned shapelets
operate on \(X_i^{\text{rel-current}}\), so they search for local patterns around
the current signal level rather than around an absolute threshold.

All current-level models are log-duration regressors. They predict
\[
  \widehat{y}_i^{\text{CLP}} = g_\vartheta(X_i^{\text{rel-current}},s_i),
\]
where \(s_i\) is the optional train-standardized scalar context. Neural
shapelet models are trained with a smooth-\(L_1\) loss in log-duration space,
\[
  \ell(e_i)
  =
  \begin{cases}
    \frac{1}{2}e_i^2, & |e_i|<1,\\
    |e_i|-\frac{1}{2}, & |e_i|\ge 1,
  \end{cases}
  \qquad
  e_i = \widehat{y}_i^{\text{CLP}}-y_i^{\text{CLP}} .
\]
This reduces the influence of extreme duration errors while preserving a
quadratic region around small errors. XGBoost uses a squared-error objective on
the same \(y_i^{\text{CLP}}\) target.

When scalar context is enabled, each neural model receives a one-channel
sequence
\(X_i^{\text{rel-current}}\in\mathbb{R}^{L}\) and the standardized
12-dimensional vector \(s_i^{\text{CLP}}\). The scalar vector is passed through
a learned affine--ReLU--dropout embedding before concatenation with the
shapelet representation. With \(L=30\), the XGBoost baseline receives a
42-column table: 30 current-relative lags followed by the same 12 scalar
features.

\begin{table}[htbp]
  \centering
  \small
  \caption{Current-level persistence models.}
  \label{tab:current-level-models}
  \renewcommand{\arraystretch}{1.25}
  \setlength{\tabcolsep}{3pt}
  \begin{tabular}{@{}>{\raggedright\arraybackslash}p{0.24\textwidth}
    >{\raggedright\arraybackslash}p{0.34\textwidth}
    >{\raggedright\arraybackslash}p{0.32\textwidth}@{}}
    \toprule
    Model & Representation & Purpose \\
    \midrule
    Shapelet MLP \cite{grabocka2014learning}
      & Minimum multiscale shapelet distances, optionally joined with scalar
        context.
      & Tests whether best-match local patterns are sufficient for duration
        regression. \\
    Shapelet Transformer \cite{grabocka2014learning,vaswani2017attention}
      & The same minimum-distance vector interpreted as pattern-response
        tokens.
      & Tests interactions among learned pattern responses before regression. \\
    Shapelet convolution \cite{grabocka2014learning}
      & Full multiscale shapelet response maps over temporal positions.
      & Preserves where a pattern occurs inside the context window. \\
    XGBoost regressor \cite{chen2016xgboost}
      & Flattened current-relative lags plus scalar context.
      & Provides a non-neural tabular baseline for the same log-duration target. \\
    \bottomrule
  \end{tabular}
\end{table}

All predictions are evaluated after conversion back to seconds through
\(\exp(\widehat{y}_i^{\text{CLP}})-1\).

\subsection{Switch Conversion}

The current-level models output a duration estimate, not a switch. Let
\(T_{\min}\) be the operational minimum duration required to justify switching.
The raw current-level switch is

\[
  r_i^{\text{CLP}}
  =
  \mathbb{1}
  \left[
    S_i \ge \theta
    \ \land\
    \widehat{D}_i^{\text{CLP}} \ge T_{\min}
  \right].
\]

In the experiments, \(\theta=10.0\) dB and \(T_{\min}=300\) seconds. The signal
gate \(S_i\ge\theta\) prevents a long predicted persistence around a
non-critical level from triggering a backup-link switch. The duration condition
then asks whether the current level is expected to remain persistent for at
least the operational switch horizon. The shared post-processing policy is
applied after this raw switch is produced.

\section{Long-Fade Detection}
\label{sec:long-fade-formulation}
Long-fade detection is a binary event-classification formulation. It is closely
related to current-level persistence because both approaches ask whether the
current fade condition is likely to persist long enough to justify an
operational switch. The difference is the target. Current-level persistence
predicts a continuous duration relative to the current signal level
\(S_i-\delta\). Long-fade detection instead predicts whether the grouped fade
event to which the current timestamp belongs is a long event.

\subsection{Task Definition}
Let \(\theta\) be the fade threshold and let \(T_{\min}\) be the minimum event
duration used to define an operationally relevant long fade. In the experiments,
\(\theta=10.0\) dB and \(T_{\min}=300\) seconds. A threshold crossing occurs at
index \(i\) when

\[
  S_i \ge \theta .
\]

The long-fade task does not define an event as a single uninterrupted run of
above-threshold samples. Such a definition would split one physical fade into
several fragments whenever the signal temporarily oscillates around the
threshold. Instead, threshold crossings are grouped with the same operational
convention used in the dataset builder. A grouped event starts at the first
crossing. Further crossings whose timestamps are within three hours from that
event start remain in the same grouped event. The event span extends from the
first to the last threshold-crossing sample in the group. Samples below
threshold inside this span are kept as part of the grouped event, so temporary
recoveries do not create separate labels.

For a grouped event \(e\), let \(a_e\) and \(b_e\) denote the first and last
crossing indices in the group. Its duration is represented on the canonical
sample grid as

\[
  D_e = (b_e-a_e+1)\Delta t,
\]

where \(\Delta t\) is the sampling interval. The event-level long-fade label is

\[
  y_e^{\text{LFD}}
  =
  \mathbb{1}[D_e \ge T_{\min}].
\]

Every valid decision timestamp \(i\) inside the grouped event span receives the
same label:

\[
  y_i^{\text{LFD}} = y_e^{\text{LFD}},
  \qquad i \in [a_e,b_e].
\]

This design makes the task a grouped-event classifier rather than a
pointwise-threshold classifier. The model is not asked whether the signal is
above threshold at the next sample; it is asked whether the current grouped fade
belongs to an event whose total duration reaches the operational minimum.

The input at a decision index is again restricted to past and current signal
values. The main sequence representation is relative to the threshold:

\[
  X_i^{\text{rel-threshold}}
  =
  [S_{i-L+1}-\theta,\ldots,S_i-\theta].
\]

This representation centers the input around the operational decision boundary:
positive values are above the fade threshold and negative values are below it.
The model may also receive scalar context features computed only from the
available past window and from online event-state information, such as the
current signal, recent slopes, local window statistics, elapsed time since the
first crossing, and the number of observed samples since the grouped event
started. Quantities that require knowledge of the future event end, such as the
final event duration or the final position fraction, are not allowed as model
inputs.

The exact long-fade scalar vector is
\[
\begin{aligned}
  s_i^{\text{LFD}} = [
  &S_i,\ S_i-\theta,\ \bar{S}_{i,L},\
  \operatorname{sd}_{i,L},\
  \min(X_i),\ \max(X_i),\
  \max(X_i)-\min(X_i),\\
  &\Delta_i^{(2)},\ \Delta_i^{(3)},\ \Delta_i^{(5)},\
  \bar{S}_{i,5},\ \operatorname{sd}_{i,5},\
  E_i^{\text{fade}},\ P_i^{\text{fade}}
  ] ,
\end{aligned}
\]
where \(E_i^{\text{fade}}\) is the elapsed time in seconds since the first
crossing of the grouped event, and \(P_i^{\text{fade}}\) is the number of
observed samples since that event start. Both are online quantities because they
depend only on the already observed beginning of the grouped event and the
current timestamp. The tabular XGBoost input \texttt{X\_lag\_scalar} is the
concatenation
\[
  [
  S_{i-L+1}-\theta,\ldots,S_i-\theta,\ s_i^{\text{LFD}}
  ],
\]
with the lag columns named from \(\operatorname{lag}_{L-1}\) to
\(\operatorname{lag}_{0}\). In the main \(L=30\) setup this gives 30
threshold-relative lags plus 14 scalar features. The TCN and
shapelet-convolution classifiers receive the threshold-relative sequence as a
sequence tensor and the same standardized scalar vector through a separate
scalar branch.

\subsection{Tested Models}
All long-fade models output a probability

\[
  \widehat{p}_i^{\text{LFD}}
  =
  P(y_i^{\text{LFD}}=1 \mid X_i),
\]
or a logit transformed into this probability. Neural classifiers are trained
with weighted binary cross-entropy with logits,
\[
  \mathcal{L}_{\mathrm{LFD}}
  =
  -\frac{1}{N}
  \sum_i
  \left[
    \alpha y_i^{\text{LFD}}\log \widehat{p}_i^{\text{LFD}}
    +
    (1-y_i^{\text{LFD}})
    \log(1-\widehat{p}_i^{\text{LFD}})
  \right],
\]
where \(\alpha\) is computed from the training split as the ratio between
negative and positive samples, with a safe fallback when no positive sample is
available. XGBoost uses the corresponding binary logistic boosted-tree
objective and handles class imbalance through the positive-class scale
parameter.

\begin{table}[htbp]
  \centering
  \small
  \caption{Long-fade detection classifiers.}
  \label{tab:long-fade-models}
  \renewcommand{\arraystretch}{1.25}
  \setlength{\tabcolsep}{3pt}
  \begin{tabular}{@{}>{\raggedright\arraybackslash}p{0.24\textwidth}
    >{\raggedright\arraybackslash}p{0.34\textwidth}
    >{\raggedright\arraybackslash}p{0.32\textwidth}@{}}
    \toprule
    Model & Input representation & Purpose \\
    \midrule
    XGBoost lag-and-scalar classifier
      & Threshold-relative lags plus train-standardized online scalar context.
      & Tests whether fixed lag features and simple signal summaries are
        sufficient for long-event classification. \\
    Temporal convolution classifier
      & Threshold-relative sequence plus scalar-context embedding.
      & Uses dilated residual convolutions to combine local changes with wider
        temporal context. \\
    Multiscale shapelet-convolution classifier
      & Multiscale distance-to-shapelet response maps plus scalar context.
      & Tests whether explicit local pattern matching improves over generic
        convolutional sequence features. \\
    \bottomrule
  \end{tabular}
\end{table}

All three models are causal at decision time because their inputs contain only
the past context, the current signal, and online event-state quantities already
observable at \(t_i\).

\subsection{Switch Conversion}
The long-fade classifier output is converted into a switch by thresholding the
predicted probability and gating the decision with the current signal level. Let
\(\pi_{\min}\) be the configured probability threshold. The raw switch is

\[
  r_i^{\text{LFD}}
  =
  \mathbb{1}
  \left[
    S_i \ge \theta
    \ \land\
    \widehat{p}_i^{\text{LFD}} \ge \pi_{\min}
  \right].
\]

The signal gate is important because the model can produce a long-fade
probability for timestamps inside the grouped event span even when the signal is
temporarily below threshold. In those cases, the probability may still be
informative about the event class, but the operational switch is activated only
when the current signal is in the fade region. The default probability threshold
used in the experiments is \(\pi_{\min}=0.5\). After this raw decision is
formed, the same shared switch post-processing policy is applied before
comparison with the Perfect Switch.

\section{Survival Persistence}
\label{sec:survival-formulation}
Survival persistence models the remaining duration of the current fade event as
a time-to-event problem. It is related to current-level persistence because both
tasks estimate a future duration, but the definition of the target is different.
Current-level persistence asks how long the signal will remain above
\(S_i-\delta\). Survival persistence asks how long the current fade event will
remain active before a stable recovery is observed.

The task is therefore naturally expressed with survival analysis notation. For
a decision time \(t_i\), let \(R_i\) be the residual duration of the current fade
event. The desired model output is a conditional survival function

\[
  S(u \mid X_i) = P(R_i > u \mid X_i),
\]

where \(u\) is a future time horizon. This output is more informative than a
single duration estimate: it can answer questions such as whether the current
fade is likely to persist for more than 60 seconds, 300 seconds, or 900 seconds.
The formulation also supports right-censored observations, a central feature of
survival data \cite{keiding2015survival}. In this project, censoring occurs when
the acquisition segment ends before stable recovery from the fade can be
observed.

\subsection{Task Definition}
Survival-persistence samples are generated only inside detected fade events.
Events are constructed within continuous acquisition segments using a
state-machine definition. The event starts when the activation condition is
satisfied,

\[
  S_i \ge \theta_{\text{on}},
\]

and remains active until a stable recovery is confirmed. A candidate recovery
starts when the recovery condition is satisfied,

\[
  S_i < \theta_{\text{off}}.
\]

The candidate becomes the event end only if the following recovery window is
stable enough: in the implemented configuration, the recovery window is
300 seconds long, contains at least a minimum number of observations, and has at
least 80\% of its samples satisfying the recovery condition. If a new activation
crossing occurs during the candidate recovery interval, the candidate is
aborted and the event remains active. This rule prevents short oscillations
around the threshold from breaking one physical fade into several artificial
events.

For an observed event \(e\), let \(a_e\) be the event start index and let
\(c_e\) be the first timestamp of the stable recovery interval. The survival
event end is \(t_{c_e}\). The confirmation time may occur later, because future
samples are needed to verify that the recovery interval is stable, but this
confirmation delay is not added to the physical target duration. For a sample
at index \(i\) inside the event,

\[
  R_i = t_{c_e} - t_i .
\]

Samples at the event end itself are excluded so that all observed survival
times remain strictly positive.

If the acquisition segment ends before a stable recovery can be confirmed, the
sample is right-censored. In that case the exact residual duration is unknown,
but it is known to be longer than the remaining observed follow-up in the
segment. Let \(t_{\text{seg,end}}\) be the last timestamp of the acquisition
segment. The stored survival label is then

\[
  R_i > t_{\text{seg,end}} - t_i .
\]

The canonical dataset represents this through lower and upper bounds:

\[
  y_{i,\text{lower}} =
  \begin{cases}
    R_i, & \text{if the event end is observed},\\
    t_{\text{seg,end}}-t_i, & \text{if right-censored},
  \end{cases}
\]

\[
  y_{i,\text{upper}} =
  \begin{cases}
    R_i, & \text{if the event end is observed},\\
    \infty, & \text{if right-censored}.
  \end{cases}
\]

Observed samples therefore have a degenerate interval
\([R_i,R_i]\), while right-censored samples have an interval
\([t_{\text{seg,end}}-t_i,\infty)\). Censored samples are kept in the dataset;
they are not dropped and they are not converted to artificial finite targets.
This is a key methodological difference from ordinary duration regression.

The saved inputs are model-independent. For every survival sample, the dataset
stores three sequence representations:
\[
  X_{\text{raw},i}=[S_{i-L+1},\ldots,S_i],
\]
\[
  X_{\text{rel-threshold},i}
  =
  [S_{i-L+1}-\theta_{\text{on}},\ldots,S_i-\theta_{\text{on}}],
\]
\[
  X_{\text{rel-current},i}
  =
  [S_{i-L+1}-S_i,\ldots,0].
\]
It also stores a train-standardized scalar context vector. In the implemented
dataset, this vector is ordered as
\[
\begin{aligned}
  s_i^{\text{SP}} = [
  &S_i,\ S_i-\theta_{\text{on}},\
  \bar{S}_{i,L},\ \operatorname{sd}_{i,L},\
  \min(X_i),\ \max(X_i),\
  \max(X_i)-\min(X_i),\\
  &\Delta_i^{(2)},\ \Delta_i^{(3)},\ \Delta_i^{(5)},\
  \bar{S}_{i,5},\ \operatorname{sd}_{i,5},\
  E_i^{\text{event}},\ P_i^{\text{event}},\\
  &A_i,\ U_i,\ B_i,\ Q_i
  ] .
\end{aligned}
\]
Here \(E_i^{\text{event}}\) and \(P_i^{\text{event}}\) are the elapsed seconds
and sample count since the event start, \(A_i\) indicates whether the current
sample is above the activation threshold, \(U_i\) is the time since the last
above-threshold sample, \(B_i\) is the current run length of consecutive
below-threshold samples, and \(Q_i\) is the fraction of above-threshold samples
in the trailing event-local window of up to five samples. These variables are
online event-state variables. They do not require the future event end or the
future recovery confirmation.

Future fields such as event end time, recovery confirmation time, event
duration, remaining duration, censoring flag, final event position, or future
threshold crossings are excluded from all survival model inputs. They are used
only for target construction, censoring metadata, evaluation, or audit tables.

\subsection{Tested Models}
Two survival model baselines are considered: XGBoost with an accelerated failure
time objective and a discrete-time causal TCN. Both models use the same
canonical survival labels, including right-censored samples. They differ in how
they represent the input history and in how they parameterize the conditional
survival function.

The input choice is explicit for each baseline. XGBoost-AFT is tabular: selected
sequence representations are flattened into lag columns and optionally
concatenated with scalar context. The discrete-time TCN is sequential: it
receives one or more signal channels of length \(L\), with the standardized
scalar vector entering through a separate branch when enabled.

\paragraph{XGBoost-AFT.}
XGBoost-AFT receives a feature vector \(\phi_i^{\text{SP}}\) built from scalar
context, flattened signal history, or their concatenation. No target bound,
event-end field, censoring flag, or observed remaining time is part of this
feature vector. The model assumes that the logarithm of the residual duration
follows a location model,
\[
  \log R_i = F(\phi_i^{\text{SP}}) + s_{\mathrm{AFT}} \varepsilon_i,
\]
where \(F\) is the boosted-tree ensemble, \(s_{\mathrm{AFT}}>0\) is the
configured scale, and \(\varepsilon_i\) follows the selected AFT error
distribution. The stored lower and upper label bounds are passed directly to
XGBoost: observed samples use \([R_i,R_i]\), while right-censored samples use
\([c_i,\infty)\), where \(c_i=t_{\text{seg,end}}-t_i\) is the observed
follow-up time.

After training, the booster outputs the location prediction
\(\widehat{\mu}_i=F(\phi_i^{\text{SP}})\). The survival probability at a horizon
\(u\) is obtained from the AFT error cumulative distribution \(F_\varepsilon\):
\[
  \widehat{S}(u\mid X_i)
  =
  1
  -
  F_\varepsilon
  \left(
    \frac{\log u-\widehat{\mu}_i}{s_{\mathrm{AFT}}}
  \right).
\]
The same location prediction is also converted into a predicted median residual
duration for diagnostic interpretation. The switch pipeline uses the survival
probability at the operational horizon.

\paragraph{Discrete-time causal TCN.}
The discrete-time TCN uses the same canonical survival dataset but converts the
continuous residual-duration labels to discrete event and censoring masks inside
the model adapter. Let the time grid be defined by intervals
\((b_{k-1},b_k]\), with an optional final open-ended interval. For each sample,
the adapter constructs an event mask \(e_{i,k}\) and an at-risk mask
\(a_{i,k}\). If the event is observed in bin \(k^\star\), then
\(e_{i,k^\star}=1\) and the at-risk mask is active up to that bin. If the sample
is right-censored, no event bin is marked and the at-risk mask is active only
for the bins that are fully observed before the censoring time.

The neural model receives one of the canonical sequence representations, or a
multichannel stack of them, and optionally the 18-dimensional scalar context. A
stack of left-padded causal residual convolutions encodes the sequence. The
output head produces one hazard logit per discrete time bin,
\[
  g_{i,k}=g_{\vartheta,k}(X_i,s_i),
  \qquad
  \widehat{h}_{i,k}=\sigma(g_{i,k}).
\]
Here \(\widehat{h}_{i,k}\) estimates the conditional hazard
\[
  P(R_i\in(b_{k-1},b_k] \mid R_i>b_{k-1}, X_i).
\]
The survival curve at the bin edges is reconstructed as the cumulative product
of the predicted non-event probabilities:
\[
  \widehat{S}(b_k\mid X_i)
  =
  \prod_{r=1}^{k}(1-\widehat{h}_{i,r}).
\]

Training minimizes a masked discrete-time survival negative log-likelihood,
implemented as binary cross-entropy over the at-risk bins,
\[
  \mathcal{L}_{\mathrm{DT}}
  =
  \frac{
    \sum_i w_i
    \sum_k a_{i,k}\,
    \operatorname{BCE}(g_{i,k},e_{i,k})
  }{
    \sum_i w_i
  } .
\]
For right-censored samples, the loss contains only bins known to have been
survived. This keeps censored observations informative without converting them
into artificial observed recoveries. The stored dataset remains
model-independent because the discrete masks are derived inside the TCN adapter
rather than saved as canonical labels.

\subsection{Switch Conversion}
The survival models output a curve rather than a binary switch. The operational
conversion evaluates that curve at the switch-relevant horizon
\(u^\star=300\) seconds. Let \(\rho_{\min}\) be the configured survival
probability threshold. The raw switch is

\[
  r_i^{\text{SP}}
  =
  \mathbb{1}
  \left[
    S_i \ge \theta
    \ \land\
    \widehat{S}(u^\star \mid X_i) \ge \rho_{\min}
  \right].
\]

In words, the model recommends switching only if the current signal is in the
fade region and the predicted probability that the event will persist for at
least another five minutes is high enough. The default uncalibrated probability
threshold is \(0.5\). An additional exploratory setting with threshold \(0.65\)
was retained for analysis, but it should be interpreted as a test-set
diagnostic rather than as a validation-calibrated operating point.

After this raw survival-derived switch is produced, the shared post-processing
policy is applied and the final sequence is compared with the task-scoped
Perfect Switch. This keeps the survival task aligned with the same operational
evaluation layer used for the other formulations.

\section{Comparison Between Formulations}
The four formulations are not interchangeable supervised problems. They share
the same signal-only input policy and the same final switch-evaluation layer,
but they estimate different intermediate quantities. This matters because native
predictive accuracy and operational switch quality are not equivalent: a model
may forecast or classify well and still switch poorly near the decision
boundary.

\Cref{tab:formulation-comparison} summarizes the question addressed by each
formulation, together with its main advantage and limitation.

\begin{table}[htbp]
  \centering
  \small
  \caption{Conceptual comparison of the four formulations.}
  \label{tab:formulation-comparison}
  \renewcommand{\arraystretch}{1.25}
  \setlength{\tabcolsep}{3pt}
  \begin{tabular}{@{}>{\raggedright\arraybackslash}p{0.19\textwidth}
    >{\raggedright\arraybackslash}p{0.27\textwidth}
    >{\raggedright\arraybackslash}p{0.24\textwidth}
    >{\raggedright\arraybackslash}p{0.19\textwidth}@{}}
    \toprule
    Formulation & Question answered & Main advantage & Main limitation \\
    \midrule
    Autoregressive forecasting
      & What will the next signal values be?
      & Keeps an inspectable future trajectory and leaves the threshold outside
        the training target.
      & Optimizes generic forecast error, not switch-critical error. \\
    Current-level persistence
      & How long will the current signal level persist?
      & Targets short-term persistence directly.
      & Does not model the end of the whole fade event and is sensitive to the
        current reference level. \\
    Long-fade detection
      & Is the current grouped event a long fade?
      & Aligns the target with the minimum operational duration.
      & Discards residual-duration information within the event. \\
    Survival persistence
      & What is \(P(R_i>u\mid X_i)\) for the current event?
      & Keeps censored samples and provides horizon-dependent persistence
        probabilities.
      & Depends strongly on the event-start, stable-recovery, and censoring
        policy. \\
    \bottomrule
  \end{tabular}
\end{table}

Because these questions generate decisions at different native timestamp sets,
cross-task comparison requires explicit alignment. The experimental protocol
therefore reports native task metrics and coverage-aware cross-task switch
metrics separately. The final comparison is a comparison of complete pipelines:
past signal window, task-specific output, raw switch conversion, and shared
post-processing against the same Perfect Switch reference.
